\documentclass[11pt]{article}
\usepackage[margin=1in]{geometry}
\usepackage{amsmath,amssymb,amsthm,booktabs,array,microtype}
\usepackage[colorlinks=true,linkcolor=blue,urlcolor=blue]{hyperref}
\newtheorem{theorem}{Theorem}
\newtheorem{lemma}{Lemma}
\newcommand{\R}{\mathbb R}
\newcommand{\Z}{\mathbb Z}
\newcommand{\vol}{\operatorname{vol}}
\newcommand{\MI}{\mu}
\newcommand{\one}{\mathbf 1}
\newcommand{\depth}{\operatorname{depth}}
\title{A smaller counterexample to the mixed-integer \texorpdfstring{\(1/e\)}{1/e} bound\\[3pt]
\large 52 integer variables and a Boolean family of outer fibers}
\author{PK research contribution}
\date{November 7, 2026}
\begin{document}
\maketitle
\begin{abstract}
We give a rational convex polytope whose mixed-integer halfspace depth is
strictly less than $1/(e2^{52})$ at every mixed-integer point. Its continuous
dimension is $274784+2^{61}$. The improvement comes from adding entire
Boolean families of outer fibers, while preserving the origin fiber by
explicit affine inequalities. This is about $2^{32}$ times smaller in
continuous dimension than the earlier $n=76$, $d=1372131+2^{93}$ example,
but it remains far above a billion variables. The argument below and its
exact arithmetic certificate have been independently reviewed; this new
construction has not been formalized in Lean.
\end{abstract}

\section{Statement and the two inequalities that drive the proof}
For a bounded set $S\subset\Z^n\times\R^d$ with measurable fibers, write
\[
 \MI(S)=\sum_{z\in\Z^n}\vol_d(S_z),\qquad
 \depth_S(y)=\inf_{\substack{H\text{ closed ambient halfspace}\\y\in H}}
                 \frac{\MI(S\cap H)}{\MI(S)}.
\]
The depth is defined when $0<\MI(S)<\infty$.
Halfspaces are proper: their normal is nonzero. In our example every
nonempty fiber is a full-dimensional convex body, so $\MI$ agrees with
the standard Euclidean Hausdorff measure $\mathcal H_d$ on $S$.

\begin{theorem}\label{thm:main}
There is a rational convex polytope $C$ with nonempty interior in
$\R^{52}\times\R^d$, where
\[
 d=274784+2^{61}=2305843009213968736,
\]
such that $S=C\cap(\Z^{52}\times\R^d)$ is nonempty and
\[
 \sup_{y\in S}\depth_S(y)<\frac1{e2^{52}}.
\]
Every nonempty integer fiber has dimension $d$. In particular this
also violates $2^{-52}(d/(d+1))^d$.
\end{theorem}

This addresses the specific constant $1/e$. It does not rule out a smaller
positive absolute constant in the weak conjecture.

Here is the mass bookkeeping before the geometry. We first construct a
body $C_0\subset\R^{52}\times\R^{3+h}$. Normalize volumes by the volume
$b$ of its origin fiber. The $2^{52}-1$ positive Boolean fibers all have
mass $m$. The outer fibers have total mass $E$, split equally among four
groups. Thus its total mass is $bW$, where
\begin{equation}\label{eq:W}
 W=(2^{52}-1)m+1+E=2^{52}m+\Delta,\qquad \Delta=1+E-m.
\end{equation}
Four halfspaces cover the origin fiber; each has normalized mass
$L=27/64+E/4$ and misses all positive Boolean fibers. We arrange
\begin{equation}\label{eq:twotests}
 L<\frac m e,\qquad \Delta>0.
\end{equation}
These cuts handle every point in the origin fiber. Every other integer
base is exposed, and its fiber has mass at most $m$. A common simplex
factor of dimension $\ell$ supplies a small cap at every point of those
fibers. It suffices to choose
\begin{equation}\label{eq:paddingtest}
 \frac{2^{52}m}{2\ell-1}<\Delta.
\end{equation}
Section~\ref{sec:caps} explains this last step, including isolation by a
closed ambient halfspace. No rounding of a center is used.

\section{The fibers and their integer bases}
Let $B_0\subset\R^3$ be the standard tetrahedron translated so that its
centroid is zero. Its barycentric coordinates $\lambda_0,\ldots,\lambda_3$
satisfy $\lambda_j(0)=1/4$. Put $f_j=\lambda_j-1/4$; these are linear
and $\sum_j f_j=0$. Define
\[
 T_a=\{x\in B_0:\lambda_j(x)\le1/4\ (j\ne a)\},\qquad
 \theta=\frac{\vol(T_a)}{\vol(B_0)}=\frac3{32},\qquad \beta=\frac{27}{64}.
\]
Indeed, the three coordinates $\lambda_j$, $j\ne a$, identify $T_a$
with $[0,1/4]^3$ and $B_0$ with a simplex of volume $1/6$; the Jacobian
cancels in the ratio. Also $\{f_j\ge0\}\cap B_0$ is a tetrahedron of
linear scale $3/4$, hence volume fraction $\beta$.

Partition the 52 integer coordinates into four blocks $I_0,\ldots,I_3$
of size $N=13$. For nonempty $I\subset\{1,\ldots,52\}$, let $a(I)$ be
the block of its least member. The intended integer support is
\begin{equation}\label{eq:support}
 Z=\{0,1\}^{52}\ \cup\
 \{-\one_J:\varnothing\ne J\subset I_a,\quad a=0,1,2,3\}.
\end{equation}
The added bases include all nonempty subsets of a block, rather than only
the negative unit vectors. There are $4(2^{13}-1)$ of them.

Choose the following rational parameters:
\begin{center}
\begin{tabular}{lll}
\toprule
$N=13$ & $h=274781$ & $\ell=2^{61}$\\
$s=500000001/10^9$ & $t=999999997/10^9$ & $q=100001/200000$\\
$\eta=399999999/400000000$ & $\varepsilon=1/10^{12}$ & \\
\bottomrule
\end{tabular}
\end{center}
Let $c_a$ have barycentric coordinates $1/4+3\varepsilon$ at $a$ and
$1/4-\varepsilon$ elsewhere. Write
\[
 F=[-1,1]^h,\quad B=B_0\times F,\quad
 \delta_s=2s-1,\quad\delta_t=1-t,\quad\delta_q=2q-1.
\]
For $k\ge1$ and $1\le r\le13$, set
\begin{align*}
 K(a,k)&=\bigl[2(c_a+s(T_a-c_a))+(k-1)c_a\bigr]\times(2qF),\\
 G(a,r)&=\bigl[t_r(T_a-c_a)-(r-1)c_a\bigr]\times(u_rF),\\
 t_r&=1-r\delta_t,\qquad u_r=1-r\delta_q.
\end{align*}
\begin{samepage}
Define $C_0$ to be the convex hull of the following assigned fibers:
\begin{center}
\begin{tabular}{lll}
\toprule
Integer base & Assigned fiber & Volume divided by $b=\vol(B)$\\
\midrule
$0$ & $B$ & $1$\\
$\one_I$, $I\ne\varnothing$ & $K(a(I),|I|)$ & $m=\theta(2s)^3(2q)^h$\\
$-\one_J$, $\varnothing\ne J\subset I_a$ & $G(a,|J|)$ &
 $\alpha_{|J|}=\theta t_{|J|}^3u_{|J|}^h$\\
\bottomrule
\end{tabular}
\end{center}
\end{samepage}
The proof must check that convexification neither creates extra integer
bases nor enlarges the assigned origin fiber. We do that explicitly next.

\section{Why these are exactly the actual fibers}
\begin{lemma}\label{lem:fibers}
The integer support of $C_0$ is exactly \eqref{eq:support}. All its nonzero
integer bases are exposed and their actual fibers are the assigned ones.
Its origin fiber is exactly $B$. The body $C_0$ is a rational polytope with
nonempty ambient interior, and all its nonempty fibers have dimension $3+h$.
\end{lemma}
\begin{proof}
Every projected coordinate lies in $[-1,1]$. If an integer projection has
coordinate $-1$ at an index in $I_a$, every generator in a convex
representation has that coordinate $-1$. Such generators all belong to
the negative Boolean block $I_a$. Thus the projection is a negative
Boolean vector in that block. If no coordinate is negative, it is Boolean.
Every base listed in \eqref{eq:support} was included, proving the census.

For exposure, put $\rho=1/53$. The positive base $\one_I$ uniquely
maximizes the functional with coefficients $1$ on $I$ and $-\rho$
elsewhere: its value is $|I|\ge1$, whereas every negative base has value
at most $13\rho<1$. Similarly $-\one_J$ uniquely maximizes the functional
with coefficients $-1$ on $J$ and $\rho$ elsewhere. Other negative bases
have smaller value, and positive bases have value at most
$\rho(52-|J|)<1\le|J|$. An exposed base can only use generators above
it in a convex representation. Its fiber is therefore exactly the assigned
convex fiber.

For the origin, define
\[
 v=\delta_s/4+(1-\delta_s)\varepsilon.
\]
Every generator $(z,x,w)$ satisfies the affine inequalities
\begin{equation}\label{eq:envelope}
 \lambda_j(x)+v\sum_i z_i\ge0\quad(j=0,\ldots,3),\qquad
 \pm w_l\le1+\delta_q\sum_i z_i\quad(l=1,\ldots,h).
\end{equation}
Here are the details. At a positive base of size $k$ and selected type $a$,
\[
 \min\lambda_j=
 \begin{cases}
 -\delta_s/4-(k-\delta_s)\varepsilon,&j\ne a,\\
 1/4+3(k-\delta_s)\varepsilon,&j=a.
 \end{cases}
\]
The first value is at least $-kv$, since the difference is
$(k-1)\delta_s(1/4-\varepsilon)\ge0$, and the second is positive.
The box bound follows from $2q=1+\delta_q\le1+k\delta_q$.

At a negative base of size $r$ in block $a$, the corresponding minima are
\[
 \min\lambda_j=
 \begin{cases}
 r(\delta_t/4+t\varepsilon),&j\ne a,\\
 1/4-3rt\varepsilon,&j=a.
 \end{cases}
\]
The first exceeds $rv$ by
$r(\delta_t-\delta_s)(1/4-\varepsilon)>0$. The second is at least $rv$
because our parameters satisfy
\begin{equation}\label{eq:geom}
 N[\delta_s+4\varepsilon(3t+1-\delta_s)]\le1.
\end{equation}
The free-box bound is precisely $u_r=1-r\delta_q$. The origin generators
also satisfy \eqref{eq:envelope}.

These inequalities persist under convex combinations. At $z=0$, they
say $x\in B_0$ and $w\in F$, so the origin fiber is contained in $B$.
The reverse inclusion holds by construction. This argument includes all
convex combinations of overlapping negative subsets.

All parameters and vertices are rational, and $t_r,u_r>0$, so all assigned
fibers are full-dimensional rational polytopes. The origin fiber spans
the vertical directions, while the positive unit bases span all integer
directions. Consequently $C_0$ has full affine span and nonempty interior.
\end{proof}

\section{The four cuts at the origin}
Put $R=159=3(52+1)$. For $j=0,\ldots,3$, define the closed halfspace
\[
 H_j=\left\{(z,x,w):f_j(x)-R\sum_{i\in I_j}z_i
                  +\eta\varepsilon\sum_{i\notin I_j}z_i\ge0\right\}.
\]
It ignores the free-box coordinates. These four halfspaces cover the origin
fiber because $\sum_j f_j=0$.

On a positive fiber of size $k$ that meets $I_j$, the continuous score is
at most $k+1$ and the integer score at most $-R+k-1$. Their sum is negative.
If its base avoids $I_j$, its selected type differs from $j$ and the maximum
total score is
\[
 [k\eta-(k-\delta_s)]\varepsilon<0,
\]
since $\eta<1-\delta_s$. Thus $H_j$ misses every positive Boolean fiber.

On a negative fiber $G(a,r)$ with $a\ne j$, the maximum score is
$r(t-\eta)\varepsilon<0$. On $G(j,r)$ the minimum score is
$rR-3rt\varepsilon>0$, so the whole fiber is included. Exactly one
negative block is captured. It follows that, with
\[
 E=4\sum_{r=1}^{13}\binom{13}{r}\alpha_r,
\]
the total volume is $bW$ from \eqref{eq:W}, and
\begin{equation}\label{eq:originmass}
 \MI\bigl((C_0\cap(\Z^{52}\times\R^{3+h}))\cap H_j\bigr)
       =b\left(\frac{27}{64}+\frac E4\right)=bL.
\end{equation}
The normals are nonzero. All omitted fibers have strict score gaps, and
the origin cap boundaries have zero volume, so this equality holds for
the closed halfspaces as stated.

\section{A small cap at every other point}\label{sec:caps}
\begin{lemma}\label{lem:padding}
For every positive integer $\ell$, put
$\Gamma_\ell=(\ell/(\ell+1))^\ell$. Then
\[
 e\Gamma_\ell<1+\frac1{2\ell-1}.
\]
Every point of an $\ell$-simplex lies in a closed cap of relative volume
at most $\Gamma_\ell$.
\end{lemma}
\begin{proof}
The exact identity
\[
 \log(e\Gamma_\ell)=\int_0^1\frac{t}{\ell+t}\,dt<\frac1{2\ell}
\]
and $e^x<1/(1-x)$ for $0<x<1$ give the first assertion. The latter
inequality follows by integrating $1/(1-t)>1$ from $0$ to $x$.
For the second, choose a largest barycentric coordinate $a$ of the point.
Then $a\ge1/(\ell+1)$. The cap where that coordinate is at least $a$
has relative volume $(1-a)^\ell\le\Gamma_\ell$, including zero volume
when $a=1$.
\end{proof}

Let $J$ be the standard rational $\ell$-simplex and set $C=C_0\times J$.
Every nonzero fiber of $C_0$ has volume at most $mb$, because
$\alpha_r\le\theta<m$. For a point of $C$ above a nonzero integer base
$z$, take the cap $A(v')\ge c$ supplied by Lemma~\ref{lem:padding} in
the $J$ coordinate. Let $g$ expose $z$ among the finitely many integer
bases. For sufficiently large finite $Q>0$, the halfspace
\[
 A(v')-c+Q\bigl(g(z')-g(z)\bigr)\ge0
\]
agrees with the cap above $z$ and misses all other fibers. Indeed, all
exposure gaps are strictly negative and $A-c$ is bounded on compact $J$.
Its normal is nonzero because the simplex-cap functional is nonzero.
Its mass is at most $mb\Gamma_\ell\vol(J)$.

For a point above zero, one of the four cuts from the previous section,
extended to ignore $J$, has mass $bL\vol(J)$. Thus every point in the
actual mixed-integer set has a containing closed halfspace of relative
mass at most
\begin{equation}\label{eq:uniform}
 \frac{\max\{L,\Gamma_\ell m\}}{W}.
\end{equation}
By \eqref{eq:twotests}, $L/W<m/(eW)<1/(e2^{52})$.
By \eqref{eq:paddingtest} and Lemma~\ref{lem:padding},
\[
 e\Gamma_\ell\,2^{52}m
 <2^{52}m+\frac{2^{52}m}{2\ell-1}
 <2^{52}m+\Delta=W.
\]
This bounds the second term in \eqref{eq:uniform} strictly below the same
target. Only the exact parameter checks remain.

\section{Exact inequalities for the displayed parameters}
All geometric hypotheses used above follow directly from the rational
parameters. The nontrivial positive margins are
\begin{align*}
 \eta-t=2(1-s)-\eta&=\frac1{2\cdot10^9},\\
 1-13\delta_t&=\frac{999999961}{10^9},&
 1-13\delta_q&=\frac{99987}{100000},\\
 1-13[\delta_s+4\varepsilon(3t+1-\delta_s)]
 &=\frac{249999993448000000143}{250000000000000000000}.
\end{align*}
In particular $0<\varepsilon<1/4$, $1/2<s,q<1$, $0<t<\eta<2(1-s)$,
$t_r,u_r>0$, and \eqref{eq:geom} hold.

The exact rational certificate described below proves
\begin{equation}\label{eq:intervals}
 \frac{1463268}{10^6}<m<\frac{1463269}{10^6},\qquad
 \frac{465669}{10^6}<E<\frac{465670}{10^6}.
\end{equation}
Put $U=1359141/500000$. The exponential series gives
\[
 e<\sum_{j=0}^8\frac1{j!}+\frac{10}{9\cdot9!}<U,
\]
because the tail after $1/9!$ has successive ratios at most $1/10$,
with later ratios strictly smaller. The margin in the second inequality
is $1717/10206000000$. Using \eqref{eq:intervals},
\begin{align*}
 \frac{1463268/10^6}{U}-\frac{27}{64}-\frac{465670/10^6}{4}
   &=\frac{40641}{2970800000}>0,\\
 \Delta&>1+\frac{465669}{10^6}-\frac{1463269}{10^6}
       =\frac3{1250}>0,\\
 \frac3{1250}-\frac{2^{52}(1463269/10^6)}{2\ell-1}
   &=\frac{139939647533841645493}{144115188075855871968750}>0.
\end{align*}
These are precisely \eqref{eq:twotests} and \eqref{eq:paddingtest}. They
complete the proof of Theorem~\ref{thm:main}. The dimension is
$d=3+h+\ell=274784+2^{61}$. Compactness, full-dimensionality, and the
fiber census were proved before adjoining $J$ and persist under this
product. All fibers have finite positive $d$-volume, so the intrinsic
Hausdorff dimension is $d$ and its measure is their finite volume sum.
Finally $\log(1+1/d)<1/d$ gives $(d/(d+1))^d>1/e$.

\subsection*{Reproducing the power certificate}
The companion files, in the same repository directory as this source, are
\begin{quote}\small\ttfamily
check\_integer\_dimension\_52.py\\
integer-dimension-52-certificate.json
\end{quote}
The program uses integer arithmetic and
\texttt{fractions.Fraction}; no floating-point estimates are premises.
It encloses the fourteen powers $(2q)^h$ and $u_r^h$, $1\le r\le13$,
at fixed denominator $D=10^{40}$.

For a nonnegative rational number $x$, initialize bounds
$\lfloor Dx\rfloor/D\le x\le\lceil Dx\rceil/D$. If integer endpoint
pairs $[L,U]$ and $[L',U']$ enclose two nonnegative quantities after
division by $D$, their product is enclosed by
\[
 \left[\frac{\lfloor LL'/D\rfloor}{D},\,
       \frac{\lceil UU'/D\rceil}{D}\right].
\]
Monotonicity of multiplication proves this assertion. Apply it at every
multiplication and squaring in binary exponentiation. For the underlying
exact values, the invariant is
$\text{accumulator}\cdot\text{base}^{r}=x^h$, where $r$ is the remaining
exponent; the stored intervals enclose these values. For an odd $r$,
multiply the accumulator by the base, then
halve the exponent after subtracting one; for an even $r$, halve it.
Square the base whenever a further iteration is needed. Induction proves
that the final accumulator interval encloses the exact power.

The JSON records the exponent and all interval endpoints at every step.
Multiplication by the positive exact factors $\theta(2s)^3$ and
$\theta t_r^3$, followed by the positive binomial sum, yields rational
bounds strictly inside \eqref{eq:intervals}. The program also verifies
every displayed scalar margin. This is a finite exact arithmetic
certificate; it is separate from the geometric proof above and from
Lean verification of the earlier $n=76$ construction.
\end{document}
